\documentclass[a4paper,11pt]{article}
\usepackage{pos}
\usepackage{lineno} % for line numbers
\linenumbers % for line numbers
%
\title{Revolutionizing the CMS High-Level Trigger System at HL-LHC with Next Generation Triggers}

\author*[a]{Bruno Alves}

\affiliation[a]{CERN, 1211 Geneva 23, Switzerland}

\emailAdd{bruno.alves@cern.ch}

\abstract{The Next Generation Triggers (NGT) project aims to research and develop new ideas and technologies for the experiment trigger systems at the HL-LHC and beyond.
  A central component of this effort is the expansion of the CMS High-Level Trigger (HLT) system towards a novel data-scouting paradigm, where all the events accepted by the Level-1 trigger are reconstructed and recorded in a high-level, storage-efficient format suitable for prompt physics analysis.
  The R$^{3}$ (Real-time Reconstruction Revolution) program addresses this challenge by fundamentally rethinking the HLT event reconstruction to deliver offline-quality performance while operating within the critical processing and storage capacity constraints of the system.
  Advanced algorithms optimized for heterogeneous computing architectures, including GPUs and distributed accelerators, combined with a reorganization of data structures, will be incorporated to significantly reduce processing time.
  With only a single opportunity to reconstruct each event, an integrated real-time calibration loop will perform on-the-fly derivation of optimal calibration constants, enabling in-situ calibration of newly acquired data.
  The output format will be strategically tuned to provide sufficient detail for comprehensive physics analyses while minimizing data volume.
  This contribution will cover the latest developments and validation studies to establish the feasibility and physics potential of the innovative NGT scouting approach in CMS.}

\FullConference{43rd International Conference on High Energy Physics (ICHEP 2026)\\
30 July  to 5 August , 2026\\
Natal, Brazil\\}

%% \tableofcontents

\begin{document}
\maketitle


\section{The NGT Concept at HLT in the CMS Detector}
%
Text
%
\section{Pixel Tracking Acceleration}
%
Text
%
\section{Results}
%
Text
%
\section{Conclusions}
%
Text
%
\begin{thebibliography}{99}
\bibitem{cms}
  The CMS Collaboration. ``The CMS Experiment at the CERN LHC''. In: JINST 3 (2008), S08004. DOI: \url{10.1088/1748-0221/3/08/S08004}.
\bibitem{ngt}
  Next Generation Triggers. (2026). NGT Evaluation Report - 2025. CERN. DOI: \url{https://doi.org/10.17181/dvqah-g0f70}
\bibitem{trackerphase2}
  The CMS Collaboration, ``The Phase-2 Upgrade of the CMS Tracker''. \url{https://cds.cern.ch/record/2272264}
\bibitem{ngtdps}
  The CMS Collaboration, ``Computing and physics performance of a high-throughput NGT Scouting stream for the HLT Phase-2 Upgrade''. \url{https://cds.cern.ch/record/2956836}
\bibitem{octdr}
  The CMS Collaboration, ``CMS Offline Software and Computing for HL-LHC, Conceptual Design Report''. \url{https://cds.cern.ch/record/2957472}
\bibitem{demonstrator}
  The CMS Collaboration, ``Next Generation Triggers Demonstrator: Time Variation of the Calibrations and System Monitoring''. \url{https://cds.cern.ch/record/2950076}
\bibitem{muonphase2}
  The CMS Collaboration, ``Muon Tracking at the CMS High-Level Trigger for the HL-LHC Upgrade''. \url{https://cds.cern.ch/record/2950077}
\bibitem{scout}
  The CMS Collaboration, ``Enriching the physics program of the CMS experiment via data scouting and data parking''. Physics Reports Volume 1115, 17 April 2025, Pages 678-772. \url{https://doi.org/10.1016/j.physrep.2024.09.006}
\bibitem{ckf}
  R. Mankel. ``A concurrent track evolution algorithm for pattern recognition in the HERA-B main tracking system''. Nucl. Instrum. Methods A395 (1997) 169. \url{https://doi.org/10.1016/S0168-9002(97)00705-5}
\bibitem{patatrack}
  A. Bocci et al. ``Heterogeneous Reconstruction of Tracks and Primary Vertices With the CMS Pixel Tracker''. Frontiers in Big Data Volume 3 - 2020 (2020). issn: 2624-909X. DOI: \url{10.3389/fdata.2020.601728}. url: ht\url{tps://arxiv.org/abs/2008.13461}.
\end{thebibliography}

\end{document}
